% GENERATED FILE. Edit canonical lessons, then run npm run generate:books.
% canonical-source-hash: fnv1a64:cfecb351cfe35098
% canonical-lessons: SA-C19-six, SA-C19-seven, SA-C19-eight, SA-C19-nine, SA-C19-ten, SA-C19-number-family

\chapter{Counting On --- Six to Ten, and What They Prove}
\label{ch:sa-numbers-6-10}
\hlchaptermodality{19}

\begin{chapteropening}
\textbf{By the end of this chapter:} \emph{I can count to ten in Sanskrit and explain why the numbers are evidence that Sanskrit, Latin and English are related.}

Complete SA-C19-number-family and demonstrate the chapter promise.
\end{chapteropening}

\section[ṣaṭ]{\sk{षट्} (ṣaṭ) — six}

\label{lesson:SA-C19-six}



\begin{quote}
You can count to five. The next five are where the family resemblance becomes impossible to miss.
\end{quote}

\subsection*{You'll want to know: \sk{षट्}}
\textbf{\sk{षट्}} (\emph{ṣaṭ}) — \textquotedblleft{}six.\textquotedblright{}

Say it beside its relatives: Latin \emph{sex}, Greek \emph{hex}, English \emph{six}. The Sanskrit form looks least like the others on the page, because \emph{ṣ} is a retroflex sound English has no letter for — but it is the same word.

This chapter is the clearest evidence a beginner can be shown that Sanskrit, Latin and English are relatives rather than neighbours. The numbers are the proof, because numbers are borrowed less often than almost anything else.

\begin{grammarlens}[title={What you've built}]
One number, and a claim you are about to see supported four more times.
\end{grammarlens}

\subsection*{Your turn}
Say these aloud:

\begin{itemize}
\raggedright
  \item \emph{ṣaṭ} — \textquotedblleft{}six\textquotedblright{}
  \item its Latin and English cousins — \emph{sex}, \emph{six}
  \item one to six, in order
\end{itemize}

\begin{tcolorbox}[breakable,colback=teal!4,colframe=teal!35!black,title={Before you move on}]
What are the Latin and English cousins of \sk{षट्}? (\emph{Sex} and \emph{six}.) Why are numbers good evidence of relatedness? (They are rarely borrowed.)
\end{tcolorbox}
\section[sapta]{\sk{सप्त} (sapta) — seven}

\label{lesson:SA-C19-seven}



\begin{quote}
The second of the five, and the resemblance is now visible without squinting.
\end{quote}

\subsection*{You'll want to know: \sk{सप्त}}
\textbf{\sk{सप्त}} (\emph{sapta}) — \textquotedblleft{}seven.\textquotedblright{}

Latin \emph{septem}. Greek \emph{hepta}. English \emph{seven}. Line them up and the shape is one shape: \emph{s-p-t} under the vowels, with Greek turning the initial \emph{s} into a rough breathing and English softening the \emph{-pt-} to \emph{-v-}.

The Sanskrit form is not the ancestor of the others — that is a common misunderstanding worth heading off. All four descend from a \textbf{shared parent nobody wrote down}. Sanskrit is simply an old, well-recorded cousin.

\begin{grammarlens}[title={What you've built}]
Two numbers, and one correction that matters: Sanskrit is a cousin of Latin and English, not their ancestor.
\end{grammarlens}

\subsection*{Your turn}
Say these aloud:

\begin{itemize}
\raggedright
  \item \emph{sapta} — \textquotedblleft{}seven\textquotedblright{}
  \item \emph{septem}, \emph{hepta}, \emph{seven} after it
  \item whether Sanskrit is the parent of English (it is not — a cousin)
\end{itemize}

\begin{tcolorbox}[breakable,colback=teal!4,colframe=teal!35!black,title={Before you move on}]
Which consonants survive across all four forms? (\emph{S}, \emph{p}, \emph{t}.) Is Sanskrit the ancestor of Latin and English? (No — a cousin; the parent was never written down.)
\end{tcolorbox}
\section[aṣṭa]{\sk{अष्ट} (aṣṭa) — eight}

\label{lesson:SA-C19-eight}



\begin{quote}
Eight — and a word that still shows its seams.
\end{quote}

\subsection*{You'll want to know: \sk{अष्ट}}
\textbf{\sk{अष्ट}} (\emph{aṣṭa}) — \textquotedblleft{}eight.\textquotedblright{}

Latin \emph{octo}, Greek \emph{okto}, English \emph{eight} (whose silent \emph{-gh-} is the ghost of a consonant that used to be pronounced).

There is an old and appealing observation about this one: the form behaves like a \textbf{dual} — the grammatical number Sanskrit uses for things that come in pairs. Eight would then be \textquotedblleft{}two fours,\textquotedblright{} counted on two hands with the thumbs left out. It is a hypothesis rather than a settled fact, and it is offered here as one, because a good story is not the same as evidence.

\begin{grammarlens}[title={What you've built}]
Three numbers, and a reminder about how to hold an etymology that is attractive but unproven: enjoy it, label it.
\end{grammarlens}

\subsection*{Your turn}
Say these aloud:

\begin{itemize}
\raggedright
  \item \emph{aṣṭa} — \textquotedblleft{}eight\textquotedblright{}
  \item \emph{octo}, \emph{eight}
  \item whether the \textquotedblleft{}two fours\textquotedblright{} story is fact or hypothesis
\end{itemize}

\begin{tcolorbox}[breakable,colback=teal!4,colframe=teal!35!black,title={Before you move on}]
Which English word is \sk{अष्ट} related to, and what is its silent \emph{-gh-}? (\emph{Eight} — the ghost of a pronounced consonant.) Is the \textquotedblleft{}two fours\textquotedblright{} reading established? (No — a hypothesis.)
\end{tcolorbox}
\section[navan]{\sk{नवन्} (navan) — nine}

\label{lesson:SA-C19-nine}



\begin{quote}
Nine — and a word you have met before meaning something else entirely.
\end{quote}

\subsection*{You'll want to know: \sk{नवन्}}
\textbf{\sk{नवन्}} (\emph{navan}) — \textquotedblleft{}nine.\textquotedblright{} Latin \emph{novem}, English \emph{nine}.

Now the catch, and it is worth meeting deliberately rather than tripping over. You already learned \textbf{\sk{नव}} (\emph{nava}) meaning \textbf{\textquotedblleft{}new.\textquotedblright{}} They look nearly identical and they are \textbf{not} the same word:

\emph{\sk{नव}} (\emph{nava}) — new. Latin \emph{novus}, English \emph{new}. \emph{\sk{नवन्}} (\emph{navan}) — nine. Latin \emph{novem}, English \emph{nine}.

Two distinct ancient words that drifted into near-identical shapes. Latin kept them apart (\emph{novus} against \emph{novem}) and so did English (\emph{new} against \emph{nine}). Sanskrit did not, quite — and a reader who assumes one word here will misread sentences.

\begin{grammarlens}[title={What you've built}]
A pair of look-alikes, separated on purpose. This is the most common way to get a word wrong: assuming that similar shape means shared origin.
\end{grammarlens}

\subsection*{Your turn}
Say these aloud:

\begin{itemize}
\raggedright
  \item \emph{navan} — \textquotedblleft{}nine\textquotedblright{}
  \item \emph{nava} — \textquotedblleft{}new\textquotedblright{}
  \item which Latin word goes with each (\emph{novem}, \emph{novus})
\end{itemize}

\begin{tcolorbox}[breakable,colback=teal!4,colframe=teal!35!black,title={Before you move on}]
What is the difference between \sk{नव} and \sk{नवन्}? (New, and nine.) Are they the same word? (No — two ancient words that drifted into similar shapes.)
\end{tcolorbox}
\section[daśa]{\sk{दश} (daśa) — ten}

\label{lesson:SA-C19-ten}



\begin{quote}
The last of them, and the one hiding inside words you use every week.
\end{quote}

\subsection*{You'll want to know: \sk{दश}}
\textbf{\sk{दश}} (\emph{daśa}) — \textquotedblleft{}ten.\textquotedblright{} Latin \emph{decem}, Greek \emph{deka}, English \emph{ten}.

The Latin form is the one that pays. \emph{Decem} gives English \textbf{decimal}, \textbf{decade}, \textbf{December} — which was the \emph{tenth} month before the calendar was rearranged and left the name stranded two months out of place.

So \sk{दश}, \emph{decem} and \emph{ten} are one word, and a fossil of the old counting survives in the name of the wrong month.

\begin{grammarlens}[title={What you've built}]
You can count to ten in Sanskrit, and every one of the second five arrived attached to a word you already knew.
\end{grammarlens}

\subsection*{Your turn}
Say these aloud:

\begin{itemize}
\raggedright
  \item \emph{daśa} — \textquotedblleft{}ten\textquotedblright{}
  \item one to ten, straight through
  \item why December is misnamed
\end{itemize}

\begin{tcolorbox}[breakable,colback=teal!4,colframe=teal!35!black,title={Before you move on}]
Which English words carry Latin \emph{decem}? (\emph{Decimal}, \emph{decade}, \emph{December}.) Why is December's name wrong? (It was the tenth month before the calendar changed.)
\end{tcolorbox}
\section[saṅkhyā-śabdāḥ]{\sk{सङ्ख्याशब्दाः} (saṅkhyā-śabdāḥ) — the numbers, and what they prove}

\label{lesson:SA-C19-number-family}



\begin{quote}
Ten numbers is enough to see something no single word could have shown you.
\end{quote}

\subsection*{You'll want to know: \sk{सङ्ख्याशब्दाः}}
Set the ten side by side with Latin and English and the correspondence is not occasional — it is regular:

\sk{एक} one · \sk{द्व} two · \sk{त्रि} three · \sk{चतुर्} four · \sk{पञ्च} five · \sk{षट्} six · \sk{सप्त} seven · \sk{अष्ट} eight · \sk{नवन्} nine · \sk{दश} ten

Where Sanskrit has \emph{p}, Latin often has \emph{p} and English often has \emph{f}. Where Sanskrit has \emph{d}, English often has \emph{t}. Those are not impressions; they are \textbf{regular correspondences}, and their regularity is what turned the study of language into something testable in the nineteenth century.

That is what the numbers prove. Not that these languages borrowed from each other, but that they inherited from something older — and that the inheritance followed rules consistent enough to predict.

\begin{grammarlens}[title={What you've built}]
You can count to ten, and you have seen the single piece of evidence that founded historical linguistics.
\end{grammarlens}

\subsection*{Your turn}
Say these aloud:

\begin{itemize}
\raggedright
  \item all ten, in order
  \item the English cousin of each
  \item what the regularity proves — inheritance, not borrowing
\end{itemize}

\begin{tcolorbox}[breakable,colback=teal!4,colframe=teal!35!black,title={Before you move on}]
What makes the correspondences evidence rather than coincidence? (They are regular.) What do they show? (Common inheritance from an older language.)
\end{tcolorbox}
